\documentclass[11pt,a4paper]{article}
\usepackage[T1]{fontenc}
\usepackage[utf8]{inputenc}
\usepackage{lmodern,amsmath,amssymb}
\usepackage[margin=25mm]{geometry}
\usepackage{longtable,booktabs,array,calc}
\usepackage[hidelinks]{hyperref}
\hypersetup{pdftitle={Why the ancients argued about the arrow and not the sling},pdfauthor={navstokgap research working notes}}
\newcommand{\tightlist}{\setlength{\itemsep}{0pt}\setlength{\parskip}{0pt}}
\setlength{\emergencystretch}{3em}
\setcounter{secnumdepth}{0}
\begin{document}
\hypertarget{why-the-ancients-argued-about-the-arrow-and-not-the-sling}{%
\section{Why the ancients argued about the arrow and not the
sling}\label{why-the-ancients-argued-about-the-arrow-and-not-the-sling}}

The ancient debate on motion is almost entirely rectilinear. Zeno's
arrow, the Vaiśeṣika arrow, Vasubandhu's denial of passage, Hui Shi's
stick and the Mohist shadow all take straight-line or static cases,
while the sling and stone, which is where Newton \textbf{defines}
centripetal force, is missing from the philosophical argument even
though slings were ordinary objects. The absence is systematic, and the
reason is that the ancient question and Newton's are at different rungs
of the same ladder: the ancients asked whether motion is possible at
all, which is a first-order question about succession, and Newton asked
what sustains a departure from straight-line motion, which is second
order and becomes askable only once inertia is assumed. Section 3 gives
the one ancient case that is genuinely circular, the whirled firebrand,
and shows it was used for a third purpose again: neither possibility nor
force, but the difference between what appears and what is. Lateral to
the main argument; no ledger promotion.

\hypertarget{newtons-definition-uses-the-sling}{%
\subsection{1. Newton's definition uses the
sling}\label{newtons-definition-uses-the-sling}}

Definition V of the \emph{Principia} introduces the concept with the
stone, in the Motte translation held in this repository:

\begin{quote}
A stone, whirled about in a sling, endeavours to recede from the hand
that turns it; and by that endeavour, distends the sling, and that with
so much the greater force, as it is revolved with the greater velocity,
and as soon as ever it is let go, flies away. That force which opposes
itself to this endeavour, and by which the sling perpetually draws back
the stone towards the hand, and retains it in its orbit, because it is
directed to the hand as the centre of the orbit, I call the centripetal
force.
\end{quote}

The sentence before it names the planets, ``perpetually drawn aside from
the rectilinear motions, which otherwise they would pursue''. So the
sling is the laboratory model of Proposition I, and the clause ``which
otherwise they would pursue'' is the load-bearing one: the straight line
is what happens by default, and only the deviation needs a cause.

\hypertarget{the-ladder-and-why-the-ancients-stood-on-its-first-rung}{%
\subsection{2. The ladder, and why the ancients stood on its first
rung}\label{the-ladder-and-why-the-ancients-stood-on-its-first-rung}}

Theorem 5 of the \href{planck-gap-paper.md}{paper} makes the difference
exact. For a signal that is an \(n\)-th order departure, with the lower
orders unknown and therefore nuisance, the quantity bounded below is
\(m\,(P^{(n)})^2\,\tau^{2n-1}\gtrsim\kappa\):

\begin{longtable}[]{@{}
  >{\raggedright\arraybackslash}p{(\columnwidth - 6\tabcolsep) * \real{0.2500}}
  >{\raggedright\arraybackslash}p{(\columnwidth - 6\tabcolsep) * \real{0.2500}}
  >{\raggedright\arraybackslash}p{(\columnwidth - 6\tabcolsep) * \real{0.2500}}
  >{\raggedright\arraybackslash}p{(\columnwidth - 6\tabcolsep) * \real{0.2500}}@{}}
\toprule\noalign{}
\begin{minipage}[b]{\linewidth}\raggedright
Rung
\end{minipage} & \begin{minipage}[b]{\linewidth}\raggedright
Question
\end{minipage} & \begin{minipage}[b]{\linewidth}\raggedright
Nuisance
\end{minipage} & \begin{minipage}[b]{\linewidth}\raggedright
Bound
\end{minipage} \\
\midrule\noalign{}
\endhead
\bottomrule\noalign{}
\endlastfoot
\(n=1\) & Is it moving at all? & position & \(mv^2\tau\ge8z^2\kappa\) \\
\(n=2\) & Is a force acting? & position and velocity &
\(F^2\tau^3/m\ge18z^2\kappa\) \\
\end{longtable}

Zeno, the Vaiśeṣika sūtras and Vasubandhu are all at \(n=1\): the thing
in dispute is whether the arrow moves, and the unknown to be beaten is
merely where it is. The sling is at \(n=2\), and so is Galileo's
comparison and all of Newton's Book I. The second rung cannot even be
\textbf{stated} without treating uniform motion as needing no account,
because otherwise curved motion is not more puzzling than straight
motion and there is no distinguished thing for a centripetal force to
explain. Ancient physics on the whole treats rest as the natural state,
so the second rung is not available to it, and the case that would
exhibit it has no work to do.

The Vaiśeṣika comes closest to opening it. VS 5.1.17 gives impulsion the
\textbf{first} motion only and hands the rest to the impression the
previous motion deposits, so the continuation of motion is already not
being charged to a continuing push. What is missing is the further step
that a \emph{straight} continuation is free while a \emph{curved} one is
not, and without that step the sling is an ordinary event rather than a
problem.

\hypertarget{the-one-circular-case-used-for-a-third-purpose}{%
\subsection{3. The one circular case, used for a third
purpose}\label{the-one-circular-case-used-for-a-third-purpose}}

Indian sources do have a whirled object: the firebrand circle,
\emph{alātacakra}, the ring of fire seen when a burning stick is swung.
Vasubandhu uses it in the
\href{../docs/classics/Abhidharmakosabhasya_motion_atoms_GRETIL.md}{Abhidharmakośabhāṣya},
Pradhan 189.23--24, arguing that because contact with the parts is
successive, the cognition of a whole is really a cognition of the parts,
\emph{avayaveṣv eva tad-buddhir alātacakravat}, ``the awareness is of
the parts only, as with the firebrand circle''. At Pradhan 33.9 the same
image appears with \emph{āśuvṛttyā}, ``by rapid action'', in asking
whether the eye grasps objects of its own size. The firebrand is
therefore an argument about \textbf{sampling}: a rapid succession
presents itself as a continuous extended thing, and the continuity is
the observer's.

That is a third use, and for this programme it is the most interesting
of the three. It is the ancient form of the point Theorem 5 makes about
records rather than about the world: below the mesh a succession is
indistinguishable from the continuous object it mimics, and above it the
succession is recoverable. The Buddhist example runs the inference in
the direction the theorem does not license, from the appearance being a
succession to the whole being unreal, but the phenomenon it names is
exactly the one being measured.

\hypertarget{consequence-for-state}{%
\subsection{4. Consequence for STATE}\label{consequence-for-state}}

The case selection in the ancient sources is explained rather than
merely noted: the arrow is the first rung and the sling the second, and
the second needs inertia before it becomes a question. The Vaiśeṣika
impression is the nearest ancient approach to the missing premise. The
firebrand supplies the one circular example and uses it for sampling,
which is the closest ancient statement of what Theorem 5 says about
records. Section 7 of the paper now carries the point in short form.

Open: whether any Greek source argues about the sling rather than merely
using it, with Ps.-Aristotle \emph{Mechanica} and Aristotle
\emph{Physics} VIII.10 on projectiles as the places to look; and whether
the \emph{alātacakra} appears in the Nyāya-Vaiśeṣika commentaries on VS
5.1.16 beside the arrow and the bird, which the Upaskāra is reported to
do and this collection cannot yet check.
\end{document}
